\chapter*{ЗАКЛЮЧЕНИЕ}
\addcontentsline{toc}{chapter}{ЗАКЛЮЧЕНИЕ}
\sloppy

Разработан интеллектуальный методический ассистент
(ИИ-методист)~--- микросервис образовательной платформы Adapstory,
поддерживающий методиста и~разработчика курса при~подготовке учебных
курсов на~основе документов заказчика с~использованием LLM,
поиска по~корпусу и~агентной оркестрации.

\textbf{Итоги выполнения задач исследования.}
Проведены 20 глубинных B2B-интервью, подтвердивших
устойчивую проблему трудоёмкой разработки корпоративных курсов.
Для~70\,\% респондентов главным ограничением оказалась
длительность разработки, 50\,\% уже прибегают к~подрядчикам
или~внешним методистам, а~60\,\% готовы рассмотреть ИИ-инструмент.
На~техническом уровне выбран и~обоснован
стек на~основе LangGraph, Dify/RAG, Knowledge~Graph,
PostgreSQL checkpointer/result store и~локального Qwen-контура;
реализованы контур обработки источников, 6-node агентный конвейер
подготовки курса, source/provenance-слой, quality seal и~snapshot
дипломных метрик H1--H3.

\textbf{Количественные результаты текущего этапа.}
Для~пилотного 19-документного прогона B2 воспроизводимо измерен
базовый слой: среднее время подготовки структуры --- 4.37~с,
среднее покрытие уровней таксономии Блума --- 61.4\,\%, доля курсов
с~покрытием $\geq 70\,\%$ --- 21.1\,\%, структурная валидность ---
100\,\%. На том же корпусе выполнен paired-прогон B2/B4:
B4 повысил среднее покрытие Блума до 83.3\,\%, дал среднюю
строковую дельту +21.9 п.п., source-grounding 100\,\% и quality seal
\texttt{PASS} в 19 из 19 запусков. Подготовлены оценочные наборы
для~следующей итерации: 30 вопросно-ответных пар для~RAG-оценки
и~100 утверждений для~факт-чекинга. Сценарная unit-экономика даёт
инфраструктурную себестоимость порядка 75~руб.\ за~курс,
LTV / CAC~=~10.0, период окупаемости 2~месяца и~точку безубыточности
на~уровне 14 клиентов.

\textbf{Статус гипотез.}
H1 технически поддержана paired-прогоном: B4 укладывается в порог
30~минут и проходит целевой уровень Блума, но продуктовая часть требует
пользовательского пилота с реальными методистами. H2 эмпирически
поддержана на 19-документном корпусе: B4 улучшает Bloom/source-grounding
относительно B2 на тех же документах. H3 предварительно поддержана
сценарной unit-экономикой: расчётная
себестоимость подготовки курса на~порядки ниже ручной разработки.

\textbf{Научный вклад.}
Предложена и~апробирована архитектура воспроизводимой
валидации LLM-систем для~поддержки проектирования учебных курсов
на~российском корпусе открытых и~официальных документов
с~зарубежным сравнительным контуром. Архитектура объединяет
методологическую рамку проектирования курса, агентную оркестрацию
с~сохраняемым состоянием
и~RAG-контур. Сформирован воспроизводимый набор артефактов
для~дальнейшей оценки: корпус документов, оценочные артефакты для~RAG
и~факт-чекинга, сценарная бизнес-модель и~протокол валидации для~пилотов.

\textbf{Практическая значимость.}
ИИ-методист представляет собой рабочий MVP
с~контуром обработки источников, подготовкой структуры курса,
ревью с~участием человека, source/provenance-слоем, quality seal
и~подготовленной основой для~контентной и~фактической валидации.
Следующие шаги~--- абляционное исследование B4, решенческие интервью,
письма о~намерениях
и~пилоты на~реальных документах заказчиков, после чего система может
быть доведена до~полного подтверждения соответствия решения выявленной
проблеме.
